\subsubsection{gprof}\label{gprof}

\href{https://sourceware.org/binutils/docs/gprof/}{gprof} je klasičan GNU alat za profajliranje programa prevedenih sa opcijom \texttt{-pg}. Kompajler i linker ubacuju dodatnu instrumentaciju koja broji pozive funkcija i omogućava da se posle izvršavanja napravi datoteka \texttt{gmon.out}.

Ovaj model je koristan kao uvod u instrumentaciono profajliranje: program se mora posebno prevesti, pokrenuti do normalnog završetka, a izveštaj se pravi naknadno. Cena merenja i sam način instrumentacije mogu promeniti ponašanje programa, pa rezultate treba čitati kao trag za dalju analizu, ne kao apsolutnu istinu.

Primer poredi namerno spor \texttt{bubble\_sort} i bibliotekarsko sortiranje:

\codefile{sort_workload.c}{c}

Prevodimo program sa \texttt{-pg}:

\begin{minted}{bash}
$ make
\end{minted}

Prvo pokretanje pravi \texttt{gmon.out}:

\begin{minted}{bash}
$ ./sort_workload
$ ls gmon.out
\end{minted}

Zatim pravimo tekstualni izveštaj:

\begin{minted}{bash}
$ gprof ./sort_workload gmon.out > gprof.txt
$ less gprof.txt
\end{minted}

\paragraph{Flat profile}\label{gprof-flat-profile}

Prvi deo izveštaja je \emph{flat profile}. On grupiše uzorke po funkciji:

\begin{minted}{text}
  %   cumulative   self              self     total
 time   seconds   seconds    calls  ms/call  ms/call  name
  ...
         ...       ...          1             bubble_sort
\end{minted}

Za ovaj primer očekujemo da se \texttt{bubble\_sort} izdvoji kao dominantna funkcija. Ako se u izveštaju ne vidi očekivana razlika, povećati veličinu ulaza ili broj ponavljanja. Kratki programi daju premalo uzoraka.

\paragraph{Call graph}\label{gprof-call-graph}

Drugi deo izveštaja je graf poziva. On povezuje vreme funkcije sa njenim pozivaocima i pozvanim funkcijama:

\begin{minted}{text}
index % time    self  children    called     name
                ...       ...       1/1          run_bubble
[1]     ...     ...       ...       1         bubble_sort
\end{minted}

Graf poziva pomaže kada ista funkcija ima više pozivalaca. Tada nije dovoljno znati da je funkcija skupa; treba znati koji tok programa je najviše koristi.

\paragraph{Ograničenja}\label{gprof-ogranicenja}

\begin{itemize}
\tightlist
\item
  Program treba da završi normalno da bi se \texttt{gmon.out} upisao.
\item
  Optimizacije, inlining i kratko izvršavanje mogu sakriti funkcije ili dati nestabilne procente.
\item
  Višedretveni programi i programi koji mnogo čekaju na I/O zahtevaju pažljivo tumačenje.
\item
  \texttt{gprof} ne zamenjuje moderne uzorkujuće alate, ali je dobar za razumevanje razlike između instrumentacije poziva i statističkog uzorkovanja.
\end{itemize}
